%% ============================================================
%%  TEMPLATE LAPORAN PRAKTIKUM - IF25-21008 JARINGAN KOMPUTER
%%  Teknik Informatika - Institut Teknologi Sumatera (ITERA)
%%  Kompilasi: XeLaTeX (2x agar nomor halaman/daftar benar)
%% ============================================================
\documentclass[12pt]{article}

% ---------- Ukuran kertas (B5 = 176 x 250 mm, sama dengan Word) ----------
% Word section2 (isi): pgMar 1440 twips = 25.4mm semua sisi, header/footer 720 = 12.7mm.
% Section1 cover 958 = 16.9mm tidak berpengaruh karena cover pakai TikZ overlay full-page.
\usepackage[paperwidth=176mm,paperheight=250mm,
            left=25.4mm,right=25.4mm,top=25.4mm,bottom=25.4mm,
            headheight=15pt,headsep=10.1mm,footskip=12.54mm]{geometry}

% ---------- Bahasa & font (XeLaTeX, font Windows asli sesuai docx) ----------
% Word: body Garamond 12pt, heading Arial Black/Bernard MT/Arial, top cover Cambria
\usepackage{fontspec}
% Opsional (jika tersedia, mis. di Overleaf): \usepackage[indonesian]{babel}
\setmainfont{Garamond}[
  BoldFont = Garamond Bold,
  ItalicFont = Garamond Italic,
  BoldItalicFont = Garamond Bold,
  BoldItalicFeatures = {FakeSlant=0.2},
]
\newfontfamily\arialblack{Arial Black}
\newfontfamily\arialfont{Arial}[
  BoldFont = Arial Bold,
  ItalicFont = Arial Italic,
  BoldItalicFont = Arial Bold Italic,
]
\newfontfamily\bernard{Bernard MT Condensed}[AutoFakeBold=1.5]
\newfontfamily\cambriafont{Cambria}
\usepackage{microtype}

% ---------- Paket pendukung ----------
\usepackage{graphicx}
\usepackage[table,dvipsnames]{xcolor}
\usepackage{tikz}
\usepackage[most]{tcolorbox}
\usepackage{fancyhdr}
\usepackage{titlesec}
\usepackage{enumitem}
\usepackage{caption}
\usepackage{array}
\usepackage{amssymb}
\usepackage{ragged2e}
\usepackage{float}

% ---------- Warna ----------
\definecolor{abuabu}{HTML}{D6D6D6}    % kotak judul sampul
\definecolor{abubanner}{HTML}{EAEAEA}  % banner judul sampul
\definecolor{abujudul}{HTML}{959595}   % teks & garis halaman judul dalam
\definecolor{abugelap}{HTML}{808080}
\definecolor{tabhead}{HTML}{4F81BD}   % header tabel
\definecolor{tabband}{HTML}{DBE5F1}   % baris selang-seling
\definecolor{tabline}{HTML}{95B3D7}   % garis tabel

% ---------- Data yang dapat diisi ----------
% Isi data peserta utama di file ini.
\newcommand{\namapeserta}{Nama Peserta Praktikum}
\newcommand{\nim}{NIM}
% Isi kelas praktikum di file ini.
\newcommand{\kelas}{RA/RB/RC/RD/RE/RF/RG/RH}
% Isi nama asisten praktikum di file ini.
\newcommand{\asistenA}{Nama Asisten Praktikum 1}
\newcommand{\asistenB}{Nama Asisten Praktikum 2}
% Isi anggota kelompok. Gunakan \\ untuk berpindah baris.
% Laporan individu: kosongkan, yaitu \newcommand{\daftaranggota}{}, agar blok anggota tidak tampil.
\newcommand{\daftaranggota}{%
  Nama Mahasiswa (NIM)\\
  Nama Mahasiswa (NIM)\\
  Nama Mahasiswa (NIM)\\
  Nama Mahasiswa (NIM)\\
  Nama Mahasiswa (NIM)\\
  Nama Mahasiswa (NIM)\\
  Nama Mahasiswa (NIM)\\
  Nama Mahasiswa (NIM)\\
  Nama Mahasiswa (NIM)}

\newcommand{\catatananggota}{Jika tugas individu cukup tulis nama dan nim saja}
% Isi informasi umum laporan praktikum di file ini.
\newcommand{\pertemuan}{1}
\newcommand{\judul}{Judul Praktikum}
\newcommand{\bulantahun}{September 2026}

% ---------- Paragraf & spasi (1,5 spasi, rata kiri-kanan) ----------
\linespread{1.4}\selectfont
\setlength{\parindent}{0pt}
\setlength{\parskip}{0pt}
\setlength{\emergencystretch}{3em}

% ---------- Header & footer (Word: Arial Black 7pt header, 8pt footer) ----------
% Garis header Word 151.7mm centered (melebar 13mm tiap sisi) -> hfoffset benar untuk GARIS.
% Tapi TEKS header/footer Word selebar badan teks -> bungkus \makebox[\textwidth] agar rata body.
\newcommand{\headtext}{%
  \arialblack
  {\fontsize{\headkirifs}{10}\selectfont\addfontfeatures{LetterSpace=\headls}LAPORAN PRAKTIKUM}\hfill
  {\fontsize{\headfs}{10}\selectfont Teknik Informatika ITERA}}
\newcommand{\headfs}{8.05}
\newcommand{\headkirifs}{7}
\newcommand{\headls}{31.5}
% Garis header Word: tipis di atas, tebal di bawah
\newcommand{\garisheader}{\hbox to\headwidth{\vbox{\hrule width\headwidth height 1.2pt
    \vskip 0.6pt\hrule width\headwidth height 3pt}}\vskip -4.8pt}

\fancypagestyle{halamandua}{%  halaman judul dalam
  \fancyhf{}
  \fancyhfoffset[HL]{13.85mm}\fancyhfoffset[HR]{12.44mm}
  \fancyhead[C]{\hspace*{1.41mm}\raisebox{-2.3mm}[0pt][0pt]{\makebox[\textwidth][s]{\headtext}}}
  \renewcommand{\headrulewidth}{4.8pt}
  \renewcommand{\headrule}{\garisheader}
  \fancyfoot[C]{\makebox[\textwidth][r]{\arialblack\fontsize{8}{9}\selectfont\thepage}}
  \renewcommand{\footrulewidth}{0.5pt}
}
\fancypagestyle{isi}{%         halaman isi
  \fancyhf{}
  \fancyhfoffset[HL]{13.85mm}\fancyhfoffset[HR]{12.44mm}
  \fancyhead[C]{\hspace*{1.41mm}\raisebox{-2.3mm}[0pt][0pt]{\makebox[\textwidth][s]{\headtext}}}
  \renewcommand{\headrulewidth}{4.8pt}
  \renewcommand{\headrule}{\garisheader}
  \fancyfoot[C]{\makebox[\textwidth][r]{\arialblack\fontsize{8}{9}\selectfont\thepage}}
  \renewcommand{\footrulewidth}{0.5pt}
}
\pagestyle{isi}

% ---------- Judul bagian (Word: Dokumentasi Garamond 20pt B, Alat/Tugas Garamond 20pt B) ----------
\titleformat{\section}[block]
  {\centering\bfseries\fontsize{20}{24}\selectfont}{}{0pt}{}
\titlespacing*{\section}{0pt}{36pt}{10pt}
\titleformat{\subsection}[block]
  {\bfseries\fontsize{11}{14}\selectfont\raggedright}{}{0pt}{}
\titlespacing*{\subsection}{0pt}{4pt}{0pt}
\titleformat{\subsubsection}[block]
  {\centering\bfseries\fontsize{20}{24}\selectfont}{}{0pt}{}
\titlespacing*{\subsubsection}{0pt}{12pt}{4pt}
% Referensi: Word SectionLabel Arial Black 15pt
\newcommand{\judulreferensi}[1]{%
  \par\addvspace{10pt}%
  {\arialblack\color{abugelap}\fontsize{15}{18}\selectfont #1}\par\vspace{8pt}}

% ---------- Daftar (bullet & sub-bullet) ----------
\setlist[itemize,1]{label=\textbullet,leftmargin=12.4mm,labelsep=6mm,labelwidth=0pt,
                    align=left,itemsep=0pt,topsep=0pt,parsep=0pt,partopsep=0pt}
\setlist[itemize,2]{label=\texttt{o},leftmargin=12.4mm,labelsep=6mm,labelwidth=0pt,
                    align=left,itemsep=0pt,topsep=0pt,parsep=0pt,partopsep=0pt}

% ---------- Caption (Word: Garamond 8pt, Gambar 1. / Tabel 1.) ----------
\DeclareCaptionFont{eightpt}{\fontsize{8}{10}\selectfont}
\captionsetup{font={eightpt},labelfont={eightpt},textfont={eightpt},labelsep=period,justification=centering,
              skip=3pt}
\renewcommand{\figurename}{Gambar}
\renewcommand{\tablename}{Tabel}

% ---------- Kotak rangkuman (Word Heading2b: Arial 12pt) ----------
\newtcolorbox{rangkuman}{enhanced,sharp corners,colback=white,colframe=black,
  % bingkai digambar sebagai stroke agar tidak hilang di zoom kecil
  boxrule=0pt,frame hidden,borderline={0.6pt}{0pt}{black},
  left=0.98mm,right=0.98mm,top=4.53mm,bottom=8.18mm,
  fontupper=\arialfont\fontsize{12}{14}\selectfont}
\newcommand{\titik}{\makebox[88mm]{\fontsize{8}{9}\selectfont\dotfill}}

% ---------- Pengaturan tabel (Word: Garamond 8pt) ----------
\renewcommand{\arraystretch}{1.0}
\newcommand{\thead}[1]{\textcolor{white}{\textbf{\fontsize{8}{10}\selectfont #1}}}

% Tampilkan pengganti yang rapi jika aset gambar opsional belum tersedia.
\newcommand{\gambarpengganti}[1]{%
  \IfFileExists{#1}{\includegraphics[width=95mm]{#1}}{%
    \begin{tikzpicture}
      \node[draw=abugelap,text=abugelap,minimum width=95mm,
            minimum height=50mm,align=center,font=\arialfont\small]
            {Gambar dokumentasi\\belum tersedia};
    \end{tikzpicture}}}

% Helper opsional untuk isi dokumentasi.
% \gbr{nama-file}{caption}: gambar dari aset/gambar/<nama-file>.png beserta caption
\newcommand{\gbr}[2]{%
  \begin{figure}[H]\centering
    \gambarpengganti{aset/gambar/#1.png}
    \caption{#2}
  \end{figure}}
% \topik{judul}: subjudul tebal kecil di dalam langkah
\newcommand{\topik}[1]{\par\vspace{8pt}{\bfseries #1}\par\vspace{2pt}}

\begin{document}
\pagenumbering{arabic}\setcounter{page}{0}

%% ==================================================================
%%  HALAMAN SAMPUL (halaman 0)
%% ==================================================================
\thispagestyle{empty}
{\linespread{1}\selectfont\begin{tikzpicture}[remember picture,overlay,
    shift={(current page.north west)},x=0.354mm,y=-0.354mm]
  % teks kiri atas (Word: Cambria 14pt)
  \node[anchor=north west,inner sep=0] at (49.2,45.2) {\cambriafont\fontsize{14}{16}\selectfont Laporan Praktikum};
  % kotak "Pertemuan" (Word PartTitle Arial Black 10pt, angka PartLabel 52pt)
  \fill[abuabu] (378.5,47.8) rectangle (448.6,128.3);
  \node[anchor=north,inner sep=0] at (413.5,51)
        {\arialblack\fontsize{10}{11}\selectfont Pertemuan};
  \node[anchor=north,inner sep=0,text=white] at (413.5,82)
        {\arialblack\fontsize{52}{52}\selectfont\pertemuan};
  % kode mata kuliah + garis (Word CompanyName: Garamond)
  \node[anchor=north west,inner sep=0] at (48.6,148.6)
        {\fontsize{12.3}{14}\selectfont IF25-21008: JARINGAN KOMPUTER};
  \draw[line width=0.5pt] (45.4,162.5) -- (449.4,162.5);
  \node[anchor=north west,inner sep=0] at (48.6,170)
        {\fontsize{18.5}{21}\selectfont Laporan Praktikum/Lab};
  % banner judul (Word TitleCover: Bernard MT Condensed 24pt Bold)
  \fill[abubanner] (15.1,236.7) rectangle (479.7,333.9);
  \node[anchor=center,inner sep=0,align=center,font=\bernard\bfseries\fontsize{24}{28}\selectfont] at (249,287.6)
        {\judul};
  % identitas (rata kanan, Word: Garamond 14pt Bold)
  \node[anchor=north east,inner sep=0,align=right,font=\bfseries\fontsize{14}{15.6}\selectfont] at (449.2,379.2)
        {\namapeserta\\ \kelas\\ \nim};
  \node[anchor=north east,inner sep=0,align=right,font=\bfseries\fontsize{14}{15.6}\selectfont] at (449.2,444.4)
        {\underline{Tim Asisten}\\ \asistenA\\ \asistenB};
  % institusi (Word: Garamond 14pt Bold)
  \node[anchor=north east,inner sep=0,align=right,font=\bfseries\fontsize{14}{15.4}\selectfont] at (449.2,545.8)
        {Teknik Informatika\\ Fakultas Teknologi Industri\\
         Institut Teknologi Sumatera\\ \bulantahun};
  % logo institusi (tiga logo, rata bawah)
  \node[anchor=base west,inner sep=0] at (42,649.4)
        {\raisebox{-2.4pt}{\includegraphics[width=57.6pt]{aset/gambar/logo-itera.png}}\hspace{23pt}%
         \includegraphics[width=114.3pt]{aset/gambar/logo-fti.png}\hspace{20.3pt}%
         \includegraphics[width=156.9pt]{aset/gambar/logo-if.png}};
\end{tikzpicture}}
\newpage

%% ==================================================================
%%  HALAMAN JUDUL DALAM (halaman 1)
%% ==================================================================
\thispagestyle{halamandua}
{\linespread{1}\selectfont\begin{tikzpicture}[remember picture,overlay,
    shift={(current page.north west)},x=0.354mm,y=-0.354mm]
  \node[anchor=north,inner sep=0,align=center,text=abujudul,font=\fontsize{9.4}{16.6}\selectfont\addfontfeatures{LetterSpace=14.5}] at (248.6,139.6)
    {FAKULTAS TEKNOLOGI INDUSTRI\\ ITERA 2026};
  % judul boleh dipecah dengan \\; garis bawah ikut turun mengikuti baris terakhir
  \node[name=jd,anchor=north,inner sep=0,align=center,text=abujudul,font=\arialblack\fontsize{23.5}{27}\selectfont\addfontfeatures{LetterSpace=-7}] at (251.8,206.2)
    {\judul\rule[-8pt]{0pt}{8pt}};
  \draw[abujudul,line width=0.6pt] (70,0 |- jd.south) ++(0,12.44) -- ++(356.4,0);
  % anggota kelompok (Word ReturnAddress: Garamond Bold), tumbuh ke atas
  % kosongkan \daftaranggota untuk laporan individu agar blok ini tidak tampil
  \ifx\daftaranggota\empty\else
  \node[anchor=south,inner sep=0,align=center,font=\bfseries\fontsize{9.6}{11.2}\selectfont] at (249,527)
    {Nama anggota kelompok:\\\daftaranggota
     \ifx\catatananggota\empty\else\\[6pt]\catatananggota\fi};
  \fi
  % alamat institusi (Word ReturnAddress: Garamond)
  \node[anchor=north,inner sep=0,align=center,font=\fontsize{9.6}{11.2}\selectfont] at (249,541.3)
    {ITERA Laboratory\\ Program Studi Teknik Informatika\\
     Fakultas Teknologi Industri\\ Institut Teknologi Sumatera};
  \node[anchor=north,inner sep=0,align=center,font=\fontsize{9.6}{11.2}\selectfont] at (249,596.7)
    {Jl. Terusan Ryacudu, Desa Way Hui, Kecamatan Jati Agung,\\
     Lampung Selatan -- Lampung 35365\\
     Telp. +62-721-8030188; 8030189};
\end{tikzpicture}}
\newpage
%% ==================================================================
%%  ISI LAPORAN (halaman 2 dst.)
%% ==================================================================
\pagestyle{isi}

% Setiap bagian laporan disimpan terpisah agar mudah diisi.
% Ganti setiap \titik dengan poin rangkuman praktikum.
% Word Heading2b: Arial bold italic, isi Garamond dengan bullet kotak.
\vspace*{-8.8mm}
{\centering\bfseries\fontsize{15}{18}\selectfont\judul\par}\vspace{1.4mm}
\begin{rangkuman}
  \linespread{1.0}\selectfont
  {\arialfont\bfseries\itshape\fontsize{12}{14}\selectfont Rangkuman praktikum}\par\vspace{4mm}
  {\arialfont\bfseries\itshape\fontsize{10}{12}\selectfont
   Poin-poin utama yang anda ketahui selama praktikum adalah:\par}\vspace{-1.5mm}
  \begin{itemize}[label=\raisebox{0.45ex}{\fontsize{3.2}{3.2}\selectfont$\blacksquare$},leftmargin=10.9mm,labelsep=3.85mm,labelwidth=0pt,
                  align=left,itemsep=0pt,topsep=0pt,parsep=0pt,
                  before=\normalfont\fontsize{9.6}{17.4}\selectfont]
    \item \titik
    \item \titik
    \item \titik
    \item \titik
    \item \titik
  \end{itemize}
\end{rangkuman}

\section{Dokumentasi Praktikum}

\subsection{Jelaskan Langkah-langkah yang dilakukan selama praktikum}

\begin{itemize}
  \item Langkah pertama adalah (silahkan jelaskan langkah-langkah
      selama praktikum)

      Lampirkan gambar yang dilakukan selama langkah pertama.

      \begin{figure}[H]
        \centering
        \gambarpengganti{aset/gambar/contoh-gambar.png}
        \caption{Ilustrasi dokumentasi praktikum}
      \end{figure}
  \item Langkah kedua adalah (silahkan jelaskan langkah-langkah
      selama praktikum)

      Jika ada tabel yang perlu dilampirkan berikut format table nya.

      \begin{table}[H]
        \centering
        \caption{Contoh penggunaan table}
        \linespread{1}\fontsize{7}{8}\selectfont
        \rowcolors{2}{tabband}{white}
        \setlength{\tabcolsep}{3pt}
        \renewcommand{\arraystretch}{1.25}
        \arrayrulecolor{tabline}
        \begin{tabular}{!{\color{tabline}\vrule}>{\centering\arraybackslash}p{9mm}!{\color{tabline}\vrule}p{34mm}!{\color{tabline}\vrule}p{34mm}!{\color{tabline}\vrule}p{34mm}!{\color{tabline}\vrule}}
          \hline
          \rowcolor{tabhead}
          \thead{No} & \thead{Judul kolom 1} & \thead{Judul Kolom 2} & \thead{Judul Kolom 3} \\
          \hline
          & & & \\ \hline
          & & & \\ \hline
          & & & \\ \hline
          & & & \\ \hline
          & & & \\ \hline
          & & & \\ \hline
          & & & \\ \hline
          & & & \\ \hline
        \end{tabular}
      \end{table}
  \item Langkah ketiga adalah (silahkan jelaskan langkah-langkah
      selama praktikum). Jika ada penggunaan list atau item di dalam
      penjelasan anda, silahkan gunakan format bullet and numbering berikut:
      \begin{itemize}
        \item List 1
        \item List 2
        \item List 3
      \end{itemize}
\end{itemize}
\par
\subsubsection{Alat dan Bahan}

Jelaskan alat dan bahan yang digunakan dalam praktikum. Deskripsikan alat atau
device yang anda gunakan. Jika menggunakan komputer Windows, cukup dengan
memberikan screenshoot dxdiag.
\par
\par
\subsubsection{Tugas}

Silahkan jelaskan jawaban dari tugas yang diberikan. Jika tidak ada tugas, maka
setiap praktikum anda wajib melampirkan dokumen laporan praktikum ini.
\par
\clearpage
\judulreferensi{Referensi}
{\linespread{1}\fontsize{11}{13}\selectfont
 \setlength{\leftskip}{22mm}\setlength{\parindent}{-22mm}
 A Forouzan, Behrouz. 2007. Data Communication and Networking Fourth
 Edition, McGraw-Hill Education, Singapore.\par}
\vspace{24pt}
{\linespread{1}\fontsize{11}{13}\selectfont\raggedright
 Lampirkan referensi yang digunakan pada modul praktikum, jika ada referensi
 tambahan lain yang anda gunakan dalam menjawab tugas, silahkan dilampirkan
 juga dengan format seperti di atas.\par}


\end{document}